% ============================================================
%  Arbeitsblatt: Pfeilklassen (Analytische Geometrie)
%  Kompilieren:  latexmk Arbeitsblatt_Vektoren.tex
%  Lösungen ein-/ausblenden: \loesungentrue / \loesungenfalse
%
%  Aufbau (nach handschriftlicher Überarbeitung):
%    Aufgabe 1  Pfeilklassen im Kästchengitter erkennen und zeichnen
%    Aufgabe 2  Pfeilklassen am Körper aus drei Würfeln
%    Aufgabe 3  Schachspiel: Pfeilklassen der Springerzüge
% ============================================================
\documentclass[11pt,a4paper]{article}

% --- Sprache & Schrift ---------------------------------------
\usepackage[ngerman]{babel}
\usepackage[T1]{fontenc}
\usepackage[utf8]{inputenc}
\usepackage{lmodern}
\usepackage{microtype}

% --- Layout --------------------------------------------------
\usepackage[a4paper,margin=2cm,top=2.5cm]{geometry}
\usepackage{fancyhdr}
\setlength{\headheight}{14pt}
\usepackage{lastpage}
\usepackage{enumitem}
\usepackage{needspace}
\usepackage{xcolor}
\usepackage{tcolorbox}
\tcbuselibrary{skins,breakable}
\setlength{\parindent}{0pt}
\setlength{\parskip}{0.5em}

% --- Mathematik ----------------------------------------------
\usepackage{amsmath,amssymb}
\usepackage{tikz}
\usetikzlibrary{arrows.meta,calc}
\usepackage{chessfss}   % Schachfiguren-Symbole
\usepackage{pgfplots}
\pgfplotsset{compat=1.18}

% Vektor-Kurzbefehle
\newcommand{\vect}[1]{\vec{#1}}
\newcommand{\spvek}[2]{\begin{pmatrix} #1 \\ #2 \end{pmatrix}}          % 2D
\newcommand{\spvekd}[3]{\begin{pmatrix} #1 \\ #2 \\ #3 \end{pmatrix}}   % 3D
\newcommand{\betrag}[1]{\left|#1\right|}

% --- Lösungen ein/aus ----------------------------------------
\newif\ifloesungen
\loesungenfalse   % <- auf \loesungentrue setzen für Lösungsblatt

\newcommand{\loesung}[1]{%
  \ifloesungen
    \begin{tcolorbox}[colback=green!5,colframe=green!50!black,
      title=Lösung,fonttitle=\bfseries\small,breakable]#1\end{tcolorbox}
  \fi}

% Platz zum Schreiben (nur ohne Lösungen)
\newcommand{\schreibplatz}[1]{\ifloesungen\else\vspace{#1}\fi}

% --- Aufgaben-Umgebung ---------------------------------------
% Optionales 2. Argument: Platz, der mindestens frei sein muss,
% sonst beginnt die Aufgabe auf einer neuen Seite.
\newcounter{aufgabe}
\newenvironment{aufgabe}[2][]{%
  \Needspace{#2}%
  \refstepcounter{aufgabe}%
  \par\medskip\noindent\textbf{Aufgabe \theaufgabe}%
  \if\relax\detokenize{#1}\relax\else\ \textit{(#1)}\fi\par\nopagebreak}%
  {\par}

% --- Kästen --------------------------------------------------
\newtcolorbox{definition}[1]{colback=orange!6,colframe=orange!70!black,
  title={#1},fonttitle=\bfseries}
\newtcolorbox{merke}[1][Merke]{colback=blue!5,colframe=blue!60!black,
  title=#1,fonttitle=\bfseries}

% --- Kopf- & Fußzeile ----------------------------------------
\newcommand{\fach}{Analytische Geometrie}

\pagestyle{fancy}
\fancyhf{}
\lhead{\fach}
\rhead{Name: \underline{\hspace{4cm}} \quad Datum: \underline{\hspace{2.5cm}}}
\cfoot{Seite \thepage{} von \pageref{LastPage}}
\renewcommand{\headrulewidth}{0.4pt}

% --- Zeichenhilfen -------------------------------------------
% \vekpfeil{Start}{Ende}{Name}{Farbe}{Beschriftungsposition}
\newcommand{\vekpfeil}[5]{%
  \draw[-Stealth,very thick,#4] #1 -- #2 node[midway,#5] {$\vect{#3}$};}
% Kästchengitter ohne Achsen: \kaestchen{xmax}{ymax}
\newcommand{\kaestchen}[2]{\draw[very thin,gray!50] (0,0) grid (#1,#2);}

% ============================================================
\begin{document}

% Farben für Lösungen (im Arbeitsblatt schwarz)
\ifloesungen
  \colorlet{gA}{blue}\colorlet{gB}{red}\colorlet{gC}{green!50!black}
  \colorlet{gX}{gray}
\else
  \colorlet{gA}{black}\colorlet{gB}{black}\colorlet{gC}{black}
  \colorlet{gX}{black}
\fi

\begin{center}
  {\LARGE\bfseries Arbeitsblatt}
  \ifloesungen\\[0.3em]{\large Lösungen}\fi
\end{center}

% ============================================================
%  Aufgabe 1: Pfeilklassen im Gitter
% ============================================================
\begin{aufgabe}{11cm}
  In der Abbildung folgenden Abbildung sind zehn Pfeile eingezeichnet.
  \begin{enumerate}[label=\alph*),nosep]
    \item Markiere alle Pfeile, die zur \textbf{selben Pfeilklasse} gehören,
          jeweils mit der gleichen Farbe.
    \item Wie viele verschiedene Pfeilklassen sind dargestellt?
    \item Begründe, warum $\vect{d}$ und $\vect{f}$ \textbf{nicht} zur selben
          Pfeilklasse gehören.
    \item Zeichne rechts drei unterschiedliche Pfeile ein, die zur selben
          Pfeilklasse gehören.
  \end{enumerate}

  \begin{center}
  \begin{tikzpicture}[scale=0.75]
    \kaestchen{22}{9}
    \vekpfeil{(1,1)}{(3,2)}{a}{gA}{above left}   % (2|1)
    \vekpfeil{(5,1)}{(7,2)}{b}{gA}{above left}   % (2|1)
    \vekpfeil{(9,4)}{(11,5)}{c}{gA}{above left}  % (2|1)
    \vekpfeil{(1,6)}{(1,4)}{d}{gB}{right}        % (0|-2)
    \vekpfeil{(6,7)}{(6,5)}{e}{gB}{right}        % (0|-2)
    \vekpfeil{(11,1)}{(11,3)}{f}{gX}{right}      % (0|2)  Gegenvektor zu d
    \vekpfeil{(3,4)}{(4,6)}{g}{gC}{left}         % (1|2)
    \vekpfeil{(8,1)}{(9,3)}{h}{gC}{left}         % (1|2)
    \vekpfeil{(8,7)}{(10,6)}{i}{gX}{below left}  % (2|-1)
    \vekpfeil{(5,4)}{(3,3)}{j}{gX}{below right}  % (-2|-1) Gegenvektor zu a
  \end{tikzpicture}
  \end{center}
  \schreibplatz{2.5cm}
  \loesung{
    a)
    \begin{itemize}[nosep]
      \item \textcolor{blue}{$\vect{a}$, $\vect{b}$, $\vect{c}$}
            (2 Kästchen nach rechts, 1 nach oben)
      \item \textcolor{red}{$\vect{d}$, $\vect{e}$}
            (2 Kästchen nach unten)
      \item \textcolor{green!50!black}{$\vect{g}$, $\vect{h}$}
            (1 Kästchen nach rechts, 2 nach oben)
      \item $\vect{f}$, $\vect{i}$ und $\vect{j}$ bilden jeweils eine
            eigene Pfeilklasse.
    \end{itemize}
    b) Es sind \textbf{sechs} verschiedene Pfeilklassen dargestellt.\\
    c) $\vect{d}$ und $\vect{f}$ sind gleich lang und parallel, aber
    \textbf{nicht gleich orientiert}: Sie zeigen in entgegengesetzte
    Richtungen.\\
    d) Individuelle Lösung: drei gleich lange, parallele und gleich
    orientierte Pfeile.
  }
\end{aufgabe}

% ============================================================
%  Aufgabe 2: Körper aus drei Würfeln
% ============================================================
\begin{aufgabe}{10cm}
  \begin{minipage}[t]{0.44\linewidth}
    \vspace{0pt}
    Der abgebildete Körper setzt sich aus drei gleich großen Würfeln
    zusammen.
    \begin{enumerate}[label=\alph*),nosep]
      \item Welche der eingezeichneten Pfeile gehören zur selben
            Pfeilklasse?
      \item Begründe, weshalb die Pfeile $\overrightarrow{JH}$,
            $\overrightarrow{KL}$ und $\overrightarrow{GL}$ nicht zur selben
            Pfeilklasse gehören, obwohl sie parallel zueinander sind.
    \end{enumerate}
  \end{minipage}\hfill
  \begin{minipage}[t]{0.54\linewidth}
    \vspace{0pt}
    \centering
    \begin{tikzpicture}[scale=2.1,
        x={(1cm,0cm)}, y={(0.5cm,0.48cm)}, z={(0cm,1cm)},
        pt/.style={circle,fill=black,inner sep=1.1pt},
        pf/.style={-Stealth,thick,blue!70!cyan}]
      % Eckpunkte
      \coordinate (A) at (0,0,0); \coordinate (B) at (1,0,0);
      \coordinate (C) at (1,1,0); \coordinate (D) at (2,1,0);
      \coordinate (E) at (2,2,0); \coordinate (N) at (0,2,0);
      \coordinate (J) at (0,0,1); \coordinate (K) at (1,0,1);
      \coordinate (L) at (1,1,1); \coordinate (M) at (2,1,1);
      \coordinate (F) at (2,2,1); \coordinate (G) at (1,2,1);
      \coordinate (H) at (0,2,1); \coordinate (I) at (0,1,1);
      % Flächen
      \fill[yellow!45!olive!25] (A) -- (B) -- (K) -- (J) -- cycle;
      \fill[yellow!45!olive!35] (B) -- (C) -- (L) -- (K) -- cycle;
      \fill[red!20] (C) -- (D) -- (M) -- (L) -- cycle;
      \fill[red!30] (D) -- (E) -- (F) -- (M) -- cycle;
      \fill[red!10] (L) -- (M) -- (F) -- (G) -- cycle;
      \fill[yellow!45!olive!12] (J) -- (K) -- (L) -- (G) -- (H) -- cycle;
      % verdeckte Kanten
      \draw[dashed,gray] (A) -- (N) -- (E)  (N) -- (H)
        (G) -- (1,2,0)  (I) -- (0,1,0) -- (C);
      % sichtbare Kanten
      \draw[thin] (A) -- (B) -- (C) -- (D) -- (E) -- (F) -- (G) -- (H)
        -- (I) -- (J) -- cycle;
      \draw[thin] (J) -- (K) -- (L) -- (M) -- (F)  (B) -- (K)  (C) -- (L)
        (D) -- (M)  (I) -- (L) -- (G);
      % Pfeile
      \foreach \s/\e in {J/H, K/L, G/L, J/L, I/G, C/E, I/B, G/D,
                         K/C, M/E, G/F, J/K, G/E, J/B}
        \draw[pf] (\s) -- (\e);
      % Punkte
      \foreach \p/\pos in {A/below left, B/below, C/below right,
          D/below right, E/right, J/left, K/below left, L/above left,
          M/right, F/above right, G/above, H/above, I/above left}
        \node[pt,label={[font=\small]\pos:$\p$}] at (\p) {};
    \end{tikzpicture}
  \end{minipage}
  \schreibplatz{4cm}
  \loesung{
    a) Zur selben Pfeilklasse gehören jeweils:
    $\{\overrightarrow{JL},\ \overrightarrow{IG},\ \overrightarrow{CE}\}$,\
    $\{\overrightarrow{IB},\ \overrightarrow{GD}\}$,\
    $\{\overrightarrow{KC},\ \overrightarrow{ME}\}$,\
    $\{\overrightarrow{JB},\ \overrightarrow{GE}\}$,\
    $\{\overrightarrow{JK},\ \overrightarrow{GF}\}$.\\
    b) $\overrightarrow{JH}$ ist doppelt so lang wie $\overrightarrow{KL}$
    (zwei Würfelkanten statt einer). $\overrightarrow{GL}$ ist zwar genauso
    lang wie $\overrightarrow{KL}$, zeigt aber in die entgegengesetzte
    Richtung (nicht gleich orientiert).
  }
\end{aufgabe}

% ============================================================
%  Aufgabe 3: Schachspiel
% ============================================================
\begin{aufgabe}[Schachspiel]{8cm}
  \begin{minipage}[c]{0.45\linewidth}
    \centering
    \begin{tikzpicture}[scale=0.6]
      \foreach \x in {1,...,8} \foreach \y in {1,...,8} {
        \pgfmathparse{mod(\x+\y,2)==0 ? "brown!75!black!85" : "yellow!8"}
        \fill[\pgfmathresult] (\x-1,\y-1) rectangle (\x,\y);}
      \draw (0,0) rectangle (8,8);
      \foreach \c [count=\i] in {a,...,h} {
        \node[font=\scriptsize,text height=1.2ex] at (\i-0.5,-0.35) {\c};
        \node[font=\scriptsize,text height=1.2ex] at (\i-0.5,8.35) {\c};}
      \foreach \i in {1,...,8} {
        \node[font=\scriptsize] at (-0.35,\i-0.5) {\i};
        \node[font=\scriptsize] at (8.35,\i-0.5) {\i};}
      \node[inner sep=0pt,font=\fontsize{17}{17}\selectfont] at (3.5,5.5)
        {\WhiteKnightOnWhite};
      \ifloesungen
        \foreach \dx/\dy in {1/2,2/1,2/-1,1/-2,-1/-2,-2/-1,-2/1,-1/2}
          \draw[-Stealth,very thick,blue] (3.5,5.5) -- ++(\dx,\dy);
      \fi
    \end{tikzpicture}
  \end{minipage}\hfill
  \begin{minipage}[c]{0.5\linewidth}
    Ein Springer hat von der Position $d6$ acht Zugmöglichkeiten.

    Wie viele \textbf{Pfeilklassen} benötigt man, um alle möglichen Züge
    anzugeben, wenn keine andere Figur stört?
  \end{minipage}
  \schreibplatz{3cm}
  \loesung{
    Man benötigt \textbf{acht} Pfeilklassen: Jeder der acht Züge
    (z.\,B.\ 2 Felder nach rechts und 1 nach oben, 1 nach rechts und 2 nach
    oben, \ldots) hat eine andere Richtung bzw.\ Orientierung. Mit diesen
    acht Pfeilklassen lassen sich die Springerzüge von \emph{jedem} Feld aus
    beschreiben -- egal, wo der Springer steht.
  }
\end{aufgabe}

\end{document}
